\section{Round 7 (Problem 30): a distinct-gap improvement of the Singer embedding}

\subsection{1) Round-7 Objective}
\textbf{Path (A): proof direction (strengthen the explicit lower-bound family).}
Round~6 globalised the Round~5 Singer subsequence by combining it with primes in short intervals, but the actual \emph{subsequence construction} from Round~5 still used only the crude fact that a $(q+1)$-point subset of a cycle of length $q^2+q+1$ has a gap of size at least the average, namely $q+1$.

The most promising direction is therefore to sharpen the \emph{cyclic geometry} of the Singer construction itself.
In this round we do exactly that.

The new theorem is:
\[
h\!\left(q^2-\left\lfloor \frac q2\right\rfloor+1\right)\ge q+1
\qquad\text{for every prime power }q.
\]
This strictly improves the Round~5 family
\[
h(q^2+1)\ge q+1.
\]
As a consequence,
\[
\limsup_{N\to\infty}\bigl(h(N)-\sqrt N\bigr)\ge \frac54,
\]
improving the Round~5 bound $\ge 1$.

The original Erd\H{o}s--Tur\'an problem remains open:
is it true that for every $\varepsilon>0$,
\[
h(N)=\sqrt N+O_\varepsilon(N^\varepsilon)\,?
\]

\subsection{2) Round-6 / Earlier Foundation Used}
We rely on the following previously vetted results.

\begin{itemize}
\item \textbf{(R5 Singer existence theorem)} For every prime power $q$ and
\[
v=q^2+q+1,
\]
there exists a perfect difference set
\[
D\subset \mathbb Z/v\mathbb Z
\qquad\text{with}\qquad |D|=q+1.
\]

\item \textbf{(R5 injective modular-to-integer transfer)} If $B\subset \mathbb Z$ is finite, reduction modulo $v$ is injective on $B$, and the image of $B$ is Sidon in $\mathbb Z/v\mathbb Z$, then $B$ is Sidon in $\mathbb Z$.

\item \textbf{(R5 cut-at-a-gap principle)} If a translate of $D\subset \mathbb Z/v\mathbb Z$ lies on an arc of length $L$, then choosing integer representatives on that arc and translating to start at $1$ gives an integer Sidon set inside $\{1,\dots,L+1\}$.

\item \textbf{(R6 final status)} The main Erd\H{o}s--Tur\'an question remains open; Round~6 produced global lower bounds from the Round~5 family, but did not improve the Singer subsequence itself.
\end{itemize}

We will \emph{strengthen} the Round~5 gap argument by exploiting a new structural property of perfect difference sets.

\subsection{3) New Insight / Tool (Round 7)}
\textbf{The cyclic gaps in a perfect difference set are pairwise distinct.}

Let
\[
0\le d_0<d_1<\cdots<d_q\le v-1
\]
be representatives of a perfect difference set $D\subset \mathbb Z/v\mathbb Z$, and define the cyclic gaps
\[
g_i:=d_{i+1}-d_i\quad(0\le i\le q-1),\qquad g_q:=d_0+v-d_q.
\]
Each $g_i$ is an ordered difference of two elements of $D$.
Since perfect difference sets have all nonzero ordered differences distinct, the $g_i$ must be pairwise distinct.

This single observation upgrades the Round~5 averaging bound:
instead of merely knowing that the largest gap is at least the average $(q^2+q+1)/(q+1)$, we now know it is the largest among \emph{$q+1$ distinct positive integers summing to $q^2+q+1$}.
That forces a much larger gap.

\subsection{4) Attack Plan (Round 7)}
The exact gap left after Round~6 is this:

\begin{quote}
Round~6 globalised the old Singer subsequence, but the intrinsic lower-bound subsequence itself still only came from the weak bound ``max gap $\ge$ average gap''.
\end{quote}

To improve it, we:
\begin{enumerate}
\item prove that the cyclic gaps in a perfect difference set are distinct;
\item prove a sharp extremal lemma for the largest of $q+1$ distinct positive integers with fixed sum $q^2+q+1$;
\item cut the Singer cycle at this improved gap;
\item transfer back to the integer setting using the Round~5 injective transfer lemma.
\end{enumerate}

This yields a strictly shorter ambient interval than in Round~5, and therefore a larger lower bound for the second-order term.

\subsection{5) Work (Round 7)}

\subsubsection{5.1. Distinct cyclic gaps in a perfect difference set}
\begin{lemma}[Distinct-gap lemma]\label{lem:distinct_gaps}
Let $q$ be a prime power, let $v=q^2+q+1$, and let
\[
D\subset \mathbb Z/v\mathbb Z
\]
be a perfect difference set of size $q+1$.
Choose representatives
\[
0\le d_0<d_1<\cdots<d_q\le v-1
\]
and define cyclic gaps
\[
g_i:=d_{i+1}-d_i\quad(0\le i\le q-1),\qquad g_q:=d_0+v-d_q.
\]
Then:
\begin{enumerate}
\item each $g_i$ is a positive integer;
\item
\[
\sum_{i=0}^q g_i=v;
\]
\item the integers $g_0,\dots,g_q$ are pairwise distinct.
\end{enumerate}
\end{lemma}

\begin{proof}
The first two assertions are immediate from the definition of the cyclic gaps.

For distinctness, note that for $0\le i\le q-1$ the gap $g_i$ is exactly the ordered difference
\[
d_{i+1}-d_i \pmod v,
\]
and
\[
g_q\equiv d_0-d_q \pmod v.
\]
Thus every $g_i$ is a nonzero ordered difference between two elements of $D$.

If $g_i=g_j$ for some $i\neq j$, then two distinct ordered pairs of elements of $D$ would yield the same nonzero ordered difference modulo $v$, contradicting the defining property of a perfect difference set.
Therefore the $g_i$ are pairwise distinct.
\end{proof}

\subsubsection{5.2. Extremal bound for the largest distinct gap}
\begin{lemma}[Largest distinct summand bound]\label{lem:max_distinct}
Let $m\ge 1$, and let
\[
1\le x_1< x_2<\cdots< x_m
\]
be distinct positive integers with sum
\[
x_1+\cdots+x_m=S.
\]
Then
\[
x_m\ge \left\lceil \frac{S+\binom{m}{2}}{m}\right\rceil.
\]
\end{lemma}

\begin{proof}
Since the $x_i$ are distinct and all at most $x_m$, the largest possible value of their sum for a fixed $x_m$ is obtained by taking
\[
x_m,\ x_m-1,\ \dots,\ x_m-(m-1).
\]
Hence
\[
S\le x_m+(x_m-1)+\cdots+(x_m-(m-1))
   =m x_m-\binom{m}{2}.
\]
Rearranging gives
\[
x_m\ge \frac{S+\binom{m}{2}}{m},
\]
and therefore
\[
x_m\ge \left\lceil \frac{S+\binom{m}{2}}{m}\right\rceil.
\]
\end{proof}

\begin{corollary}[Improved gap lower bound for Singer parameters]\label{cor:gap_lb}
With the notation of Lemma~\ref{lem:distinct_gaps}, let
\[
G:=\max_{0\le i\le q} g_i.
\]
Then
\[
G\ge \left\lfloor \frac{3q}{2}\right\rfloor+1.
\]
\end{corollary}

\begin{proof}
Apply Lemma~\ref{lem:max_distinct} with
\[
m=q+1,\qquad S=v=q^2+q+1.
\]
Since the $g_i$ are distinct positive integers summing to $v$, we obtain
\[
G\ge \left\lceil \frac{q^2+q+1+\binom{q+1}{2}}{q+1}\right\rceil
 = \left\lceil \frac{q^2+q+1+\frac{q(q+1)}2}{q+1}\right\rceil
 = \left\lceil \frac{3q}{2}+\frac{1}{q+1}\right\rceil.
\]
Because $1/(q+1)>0$, this equals $\lfloor 3q/2\rfloor+1$.
\end{proof}

\subsubsection{5.3. Improved Singer embedding}
\begin{theorem}[Improved Singer subsequence]\label{thm:improved_singer}
For every prime power $q$, there exists an integer Sidon set
\[
S\subseteq \left\{1,2,\dots, q^2-\left\lfloor \frac q2\right\rfloor+1\right\}
\]
with
\[
|S|=q+1.
\]
Equivalently,
\[
h\!\left(q^2-\left\lfloor \frac q2\right\rfloor+1\right)\ge q+1.
\]
\end{theorem}

\begin{proof}
Let
\[
v:=q^2+q+1,
\]
and let
\[
D\subset \mathbb Z/v\mathbb Z
\]
be a perfect difference set of size $q+1$, guaranteed by the Round~5 Singer existence theorem.

Choose representatives
\[
0\le d_0<d_1<\cdots<d_q\le v-1
\]
and define cyclic gaps $g_0,\dots,g_q$ as in Lemma~\ref{lem:distinct_gaps}.
Let
\[
G:=\max_i g_i.
\]
By Corollary~\ref{cor:gap_lb},
\[
G\ge \left\lfloor \frac{3q}{2}\right\rfloor+1.
\]

Now cut the cycle at a gap of size $G$.
That is, translate $D$ so that the point \emph{after} this largest gap becomes $0$, and then choose integer representatives on the complementary arc.
By the Round~5 cut-at-a-gap principle, we obtain a set
\[
B\subset \mathbb Z
\]
whose reduction modulo $v$ is a translate of $D$, and whose diameter satisfies
\[
\operatorname{diam}(B)\le v-G.
\]
Using the bound on $G$ gives
\[
\operatorname{diam}(B)\le q^2+q+1-\left(\left\lfloor \frac{3q}{2}\right\rfloor+1\right)
= q^2+q-\left\lfloor \frac{3q}{2}\right\rfloor.
\]
Since
\[
q^2+q-\left\lfloor \frac{3q}{2}\right\rfloor
=
q^2-\left\lfloor \frac q2\right\rfloor
< v,
\]
reduction modulo $v$ is injective on $B$.
Moreover the image of $B$ modulo $v$ is a translate of the perfect difference set $D$, hence is Sidon in $\mathbb Z/v\mathbb Z$.

By the Round~5 injective modular-to-integer transfer lemma, $B$ is Sidon in $\mathbb Z$.
Finally translate $B$ so that its minimum becomes $1$:
\[
S:=B-(\min B)+1.
\]
Then
\[
|S|=|B|=q+1
\]
and
\[
S\subseteq \{1,2,\dots,\operatorname{diam}(B)+1\}
   \subseteq \left\{1,2,\dots, q^2-\left\lfloor \frac q2\right\rfloor+1\right\}.
\]
Thus $S$ is an integer Sidon set of the required size in the claimed interval.
\end{proof}

\subsubsection{5.4. Consequence for the second-order term}
\begin{corollary}[Improved limsup]\label{cor:limsup_5_4}
One has
\[
\limsup_{N\to\infty}\bigl(h(N)-\sqrt N\bigr)\ge \frac54.
\]
\end{corollary}

\begin{proof}
Set
\[
N_q:=q^2-\left\lfloor \frac q2\right\rfloor+1,
\]
where $q$ ranges over prime powers.
By Theorem~\ref{thm:improved_singer},
\[
h(N_q)\ge q+1.
\]
Therefore
\[
h(N_q)-\sqrt{N_q}\ge q+1-\sqrt{N_q}.
\]

Now
\[
\left\lfloor \frac q2\right\rfloor\ge \frac q2-\frac12,
\]
so
\[
N_q\le q^2-\frac q2+\frac32.
\]
Hence
\[
q+1-\sqrt{N_q}\ge q+1-\sqrt{q^2-\frac q2+\frac32}.
\]
Let
\[
\alpha_q:=q-\frac14.
\]
Then
\[
\alpha_q^2=q^2-\frac q2+\frac1{16},
\]
so
\[
q^2-\frac q2+\frac32=\alpha_q^2+\frac{23}{16}.
\]
Therefore
\[
\sqrt{q^2-\frac q2+\frac32}
=\alpha_q+\frac{\frac{23}{16}}{\alpha_q+\sqrt{q^2-\frac q2+\frac32}}.
\]
As $q\to\infty$, the correction term tends to $0$, so
\[
\sqrt{q^2-\frac q2+\frac32}=q-\frac14+o(1).
\]
Consequently
\[
q+1-\sqrt{q^2-\frac q2+\frac32}=\frac54+o(1),
\]
and thus
\[
h(N_q)-\sqrt{N_q}\ge \frac54+o(1).
\]
Taking the limsup over prime powers $q$ gives
\[
\limsup_{N\to\infty}\bigl(h(N)-\sqrt N\bigr)\ge \frac54.
\]
\end{proof}

\subsection{6) Adversarial Verification}
We check the new argument carefully.

\begin{itemize}
\item \textbf{Are the cyclic gaps really ordered differences in the modular sense?}
Yes. For $0\le i\le q-1$, one has $g_i=d_{i+1}-d_i$ as an integer and hence modulo $v$.
For the wrap-around gap, $g_q=d_0+v-d_q$, so modulo $v$ this is $d_0-d_q$.
Thus every cyclic gap is indeed a nonzero ordered difference.

\item \textbf{Could two equal cyclic gaps come from the same ordered pair?}
No. The ordered pairs $(d_{i+1},d_i)$ are distinct as $i$ varies cyclically, so equality of two gaps would give two distinct ordered pairs with the same nonzero difference, contradicting perfectness.

\item \textbf{Is the extremal bound in Lemma~\ref{lem:max_distinct} sharp enough?}
Yes. It is the best possible bound obtainable from only ``distinct positive integers with fixed sum'', and it is exactly the information needed to improve the Round~5 average-gap argument.

\item \textbf{Is the new ambient length formula correct?}
From the cut-at-gap step, the final interval length is
\[
\operatorname{diam}(B)+1\le (v-G)+1.
\]
Since $G\ge \lfloor 3q/2\rfloor+1$ and $v=q^2+q+1$, this becomes
\[
v-G+1\le q^2+q+1-\left(\left\lfloor\frac{3q}{2}\right\rfloor+1\right)+1
      =q^2-\left\lfloor\frac q2\right\rfloor+1.
\]
So there is no off-by-one error.

\item \textbf{Does injectivity of reduction modulo $v$ hold?}
Yes, because
\[
\operatorname{diam}(B)\le q^2-\left\lfloor \frac q2\right\rfloor < q^2+q+1=v.
\]
Therefore no two distinct elements of $B$ can be congruent modulo $v$.

\item \textbf{Does the limsup computation really tend to $5/4$?}
Yes. The parameter
\[
N_q=q^2-\left\lfloor \frac q2\right\rfloor+1
\]
satisfies
\[
N_q=q^2-\frac q2+O(1),
\]
so
\[
\sqrt{N_q}=q-\frac14+o(1),
\]
which yields
\[
q+1-\sqrt{N_q}=\frac54+o(1).
\]
\end{itemize}

\subsection{7) Final}
\textbf{UNRESOLVED (BUT STRICTLY ADVANCED).}

Round~7 proves a genuinely stronger explicit lower-bound family than any previous round:
\[
h\!\left(q^2-\left\lfloor \frac q2\right\rfloor+1\right)\ge q+1
\qquad(q\text{ a prime power}),
\]
which strictly improves the Round~5 family $h(q^2+1)\ge q+1$.

As a consequence,
\[
\limsup_{N\to\infty}\bigl(h(N)-\sqrt N\bigr)\ge \frac54,
\]
improving the Round~5 lower bound $\ge 1$.

The original Erd\H{o}s--Tur\'an problem --- whether for every $\varepsilon>0$,
\[
h(N)=\sqrt N+O_\varepsilon(N^\varepsilon)
\]
--- remains open.

\subsection{8) Completion Estimate}
\textbf{COMPLETION: 80\%}

\subsection{9) References}
\begin{itemize}
\item T.\ F.\ Bloom (ed.), \emph{Erd\H{o}s Problem \#30}, \texttt{https://www.erdosproblems.com/30}, accessed 7 March 2026.
\item J.\ Singer, \emph{A theorem in finite projective geometry and some applications to number theory}, Trans.\ Amer.\ Math.\ Soc.\ \textbf{43} (1938), 377--385.
\end{itemize}
